% !TeX root = ../../Book3.tex
\section{指数与绕数}
\label{b3:sec:1001}

第2章用曲线收缩证明积分不变，本章再把曲线对各个孔的总绕行抽取出来。全局 Cauchy 定理最终只需要这些整数同时为零，而不要求曲线可以收缩。

\subsection{指数}
\label{b3:subsec:100101}

对分段光滑闭曲线 \(\gamma\) 及 \(a\notin\gamma\)，定义\term[quxian-zhishu]{曲线指数}[index of a curve]
\[
 \operatorname{Ind}(\gamma,a)=\frac1{2\pi\mathrm{i}}\int_\gamma\frac{\mathrm{d}z}{z-a}.
\]
沿曲线可连续提升 \(\arg(\gamma-a)\)，径向模的对数在闭路前后相同，积分只留下辐角总变化。因而指数为整数。它不是曲线某一点的局部属性，也不是闭路欧氏长度的函数。多次走同一圈使指数乘相同整数，反向使其变号。

\ProseParagraphBreak
指数作为 \(a\) 的函数在曲线外连续，整数性使它在每个补集连通分支上恒定。无界分支上它为0：当 \(|a|\) 很大时，积分由长度估计趋于0，所以这个整数最终为0。正向 Jordan 圈在内部指数1，外部0，使用的是平面分离和正向辐角变化；自交曲线则可能有多个不同整数区域。

\subsection{同伦不变性}
\label{b3:subsec:100102}

\begin{proposition}\label{b3:prop:10-index-homotopy}
在 \(\mathbb C\setminus\{a\}\) 中同伦的闭路具有相同指数。对有限整数线性组合，指数按系数相加。
\end{proposition}
\begin{proof}
可沿同伦的每条闭路连续选局部辐角，起终点差的整数在参数局部不变；紧参数区间由有限小辐角邻域覆盖，所以始终相同。若同伦路径不分段光滑，先以小盘内折线替换；圆盘内的原函数和连续辐角保证不改变指数。线性组合结论由积分线性直接得到。
\end{proof}
不经过点 \(a\) 是全部证明的条件。一个圆周连续缩成穿过 \(a\) 的点时，指数不再有定义，故不能以这样的变形宣称绕数已变为0。

\ProseParagraphBreak
指数作为补集上的局部常整数，可以在简单几何分解中直接计算。正向 Jordan 圈在内部为1、外部为0；两圈相加，在共同内部为2，在只被一圈围住处为1。若第二圈反向，共同内部为0。把两圈的公共部分画成一次边界并忽略方向，会遗漏这些权重。对自交的八字圈，两个瓣可能分别为1与 \(-1\)，取决于走向；交点本身属于曲线，指数在那里不定义。

\ProseParagraphBreak
无界补分支的指数为0也可用估计证明。曲线像落在某个有限圆盘中，当 \(|a|\) 大时，\(1/(z-a)=-a^{-1}\sum_{k\ge0}(z/a)^k\) 在曲线上一致收敛，各项全是整函数的导数，闭路积分0。再由局部常性把0传遍无界分支。这个事实保证非零指数区域不会跑到无穷远，是一般留数和只涉及有限紧填充的一个依据。

\subsection{指数与围道积分}
\label{b3:subsec:100103}

闭路积分中 \(1/(z-a)\) 专门测量绕 \(a\) 的次数。更一般有理函数的积分则是有限个这种测量的加权和，权重为留数。例如在避开0和1的区域中，\(1/(z(z-1))=-1/z+1/(z-1)\)，故
\[
 \frac1{2\pi\mathrm{i}}\int_\gamma\frac{\mathrm{d}z}{z(z-1)}
 =-\operatorname{Ind}(\gamma,0)+\operatorname{Ind}(\gamma,1).
\]
绕两个点各一次时两个贡献相消，积分为0，并不说明闭路不绕任何点。选定被积函数可能看不见某些拓扑数据；要判定全纯函数的全部周期，必须考虑整个区域及所有闭路。

\ProseParagraphBreak
指数为整数的证明可直接用连续辐角。把闭路相对于 \(a\) 归一化为 \(\eta(t)=(\gamma(t)-a)/|\gamma(t)-a|\)，有限小弧上逐段选辐角并在重叠处加 \(2\pi\) 整数以衔接，得连续 \(\theta(t)\)。闭路终点给 \(\theta(1)-\theta(0)=2\pi n\)，分段求导并积分即得 \(\operatorname{Ind}(\gamma,a)=n\)。改变最初辐角只加同一个 \(2\pi\) 整数，不改变差。这个整数解释为何随避开曲线的连续移动只能局部恒定。

\ProseParagraphBreak
对圈 \(\Gamma=\sum n_j\gamma_j\)，指数按系数相加。重走一条曲线两次给指数2，反向给负指数，完全可以出现负值。一般围道不是一个“内部”集合的正向边界；即使画出多重自交，也应逐个区域记录整数权重。Cauchy 与留数公式中的指数正是补偿这种几何复杂性的系数。
